% =====================================================================
%  CAPÍTULO NUEVO: RESULTADOS Y DISCUSIÓN
%  Insertar en main.tex ENTRE 04-desarrollo y 05-conclusiones (que pasa a
%  ser el capítulo 6). Reutiliza las tablas ya existentes de produccion.tex
%  (tab:metricas-reconocimiento, tab:umbral-reconocimiento,
%  tab:metricas-rutinas); no duplica cifras.
%  Todas las cifras provienen de ml/results/*.json o del texto vigente de
%  la memoria. Lo no medido se marca con \pendiente{}.
% =====================================================================
\chapter{RESULTADOS Y DISCUSIÓN}
\label{cap:resultados}

Este capítulo reúne los resultados del proyecto, los contrasta con los
objetivos específicos y con los antecedentes citados, y declara las amenazas a
la validez que limitan su interpretación. Las cifras de inteligencia artificial
proceden de los archivos de resultados del repositorio
(\archivo{ml/results/}); cuando una medida no se realizó se indica de forma
explícita.

\section{Resultados por objetivo específico}
\label{sec:res-objetivos}

La Tabla~\ref{tab:res-objetivos} resume, para cada objetivo específico, lo que
se construyó, la evidencia disponible y el grado de cumplimiento.

\begin{table}[H]
\centering
\caption{Cumplimiento de los objetivos específicos}
\label{tab:res-objetivos}
\small
\begin{tabularx}{\textwidth}{|L{2.9cm}|Y|L{3.5cm}|}
\hline
\rowcolor{encabezado} \textbf{Objetivo} & \textbf{Resultado} & \textbf{Cumplimiento} \\
\hline
OE1. Guiar sobre el uso correcto de máquinas &
Catálogo de 35 tipos de máquinas con guía de uso, errores comunes y
prevención de lesiones; escáner con reconocedor local (acierto top-1 de 38,5\,\%
en campo simulado, Tabla~\ref{tab:metricas-reconocimiento}), detector en la
nube con cajas (YOLO26m, mAP50 de 48,5\,\% en validación, Tabla~\ref{tab:metricas-detector-yolo26}) y ficha anclada en el entorno con ARCore. &
Parcial: guía completa; reconocimiento modesto y sin prueba en gimnasios
reales; anclaje implementado pero no validado en un teléfono. \\ \hline
OE2. Instruir en la ejecución correcta &
95 animaciones de ejecución con músculo objetivo y dificultad, asociadas a las
máquinas; entrenamiento guiado paso a paso. &
Parcial: contenido instructivo, sin verificación automática de la postura. \\ \hline
OE3. Rutinas personalizadas y estructuradas &
Generador basado en reglas con recomendador de TensorFlow Lite y respaldo por
reglas (Tabla~\ref{tab:metricas-rutinas}). &
Cumplido; la calidad de las rutinas para el usuario no se evaluó. \\ \hline
OE4. Herramientas de planificación &
Plan semanal por días disponibles y registro de la racha de entrenamiento. &
Cumplido. \\ \hline
OE5. Registro de ingesta calórica &
Búsqueda de alimentos en FoodData Central mediante una Cloud Function con
cuotas por usuario, cálculo de porciones y del IMC. &
Cumplido dentro del alcance (conteo calórico, no consejo nutricional). \\ \hline
OE6. Seguimiento del progreso físico &
Registro de peso, estatura y edad con cálculo y clasificación del IMC,
sincronizado con Firestore. & Cumplido. \\ \hline
OE7. Reportes personalizados &
Reportes de progreso y metas; reporte en PDF por período. &
Cumplido: el PDF se generó y se abrió en el dispositivo de prueba (CP-41,
Figura~\ref{fig:reporte-pdf}). \\
\hline
\end{tabularx}
\fuente{Elaboración propia a partir del capítulo~\ref{cap:desarrollo} y del
repositorio del proyecto.}
\end{table}

\section{Resultados de los modelos de inteligencia artificial}
\label{sec:res-ia}

\subsection{Reconocimiento de máquinas}

Con 399 fotografías de 35 clases, validación cruzada de cinco pliegues y 30
aumentos por fotografía, el modelo acierta la máquina en el 38,5\,\% de los
casos en el escenario de campo simulado y en el 61,4\,\% entre sus tres
primeras sugerencias, frente al 2,9\,\% que lograría el azar
(Tabla~\ref{tab:metricas-reconocimiento}). Al fijar el umbral de confianza en
0,8, la aplicación acepta automáticamente solo el 15,8\,\% de las imágenes, con
una precisión del 85,4\,\% (Tabla~\ref{tab:umbral-reconocimiento}).

El análisis por clase muestra que el rendimiento es muy desigual: de las 33
clases evaluables, 25 quedan por debajo del 50\,\% de acierto en el campo
simulado, mientras que la caminadora alcanza el 92,9\,\% y la bicicleta
estática el 66,7\,\%. Las confusiones se concentran entre máquinas de aspecto
similar: el 50\,\% de las muestras de validación de la máquina combinada de
extensión y curl de piernas se clasificó como curl femoral, y el 26,5\,\% de
las de press de hombro como press de pecho
(\archivo{matriz\_confusion\_campo} en
\archivo{recognition\_metrics.json}). Este patrón es coherente con un conjunto
de entrenamiento pequeño y con imágenes de catálogo, y sugiere que la mejora
proviene de más fotografías tomadas en gimnasios reales antes que de una
arquitectura distinta (Sandler et~al., 2018).

\subsection{Detector de máquinas en la nube}

El detector de la nube responde a una limitación distinta a la del clasificador
local: además de la clase, entrega la caja de la máquina en la imagen, que es lo
que permite anclar la ficha en el entorno. El proyecto tuvo dos versiones. La
primera (YOLO11x, entrenada en Colab con un conjunto anotado en Roboflow)
obtuvo, en una partición de prueba de 43 imágenes, una precisión del 79,8\,\%,
una exhaustividad del 29,9\,\% y una mAP50 del 39,8\,\%
(Tabla~\ref{tab:metricas-detector}). La versión final (YOLO26m, entrenada y
publicada en Ultralytics Platform con un conjunto anotado con polígonos) alcanzó
en la partición de validación una precisión del 57,7\,\%, una exhaustividad del
45,0\,\%, una mAP50 del 48,5\,\% y una mAP50--95 del 44,8\,\%
(Tabla~\ref{tab:metricas-detector-yolo26}).

Frente a la primera versión, sube la exhaustividad (de 29,9\,\% a 45,0\,\%) y la
mAP50 (de 39,8\,\% a 48,5\,\%), y baja la precisión (de 79,8\,\% a 57,7\,\%): el
detector final encuentra más máquinas, pero también señala algunas que no son
o se equivoca de clase con más frecuencia. La comparación es indicativa y no
concluyente. Los conjuntos de datos y las particiones son distintos (prueba
frente a validación) y la validación también sirve para elegir el punto de
control, lo que favorece a la versión final. Además, la pérdida de
clasificación en validación es unas diez veces la de entrenamiento, una señal
habitual de sobreajuste. Lo que sí permiten afirmar estos resultados es que el
procedimiento completo (anotación, entrenamiento, publicación del servicio y
consumo desde la aplicación) está construido y es repetible.

Tampoco son comparables con las del clasificador local, que se evaluó con
validación cruzada sobre fotografías de catálogo de 35 clases y que siempre
devuelve una máquina, mientras que el detector puede no encontrar ninguna. Con
la evidencia disponible no puede afirmarse que el detector supere al
clasificador local; esa comparación exige evaluar ambos con las mismas
fotografías de un gimnasio real (sección~\ref{sec:amenazas}). La literatura
sobre detectores de la familia YOLO reporta sus resultados sobre conjuntos de
referencia de gran tamaño, como COCO (Terven et~al., 2023); el proyecto parte de
un conjunto propio mucho más pequeño, y el rendimiento limitado que se observa
es el que cabe esperar en ese régimen.

\subsection{Recomendador de rutinas}

El recomendador reproduce las reglas con una exactitud del 99,8\,\% en datos
sintéticos y del 97,1\,\% en una cuadrícula exhaustiva de 28\,512 perfiles
(Tabla~\ref{tab:metricas-rutinas}). Dado que las etiquetas se generaron con las
mismas reglas, estas cifras miden la fidelidad del modelo y no demuestran que
sus rutinas sean mejores que las del generador por reglas, cuya base de
conocimiento se apoya en la literatura de prescripción del entrenamiento de
fuerza (Kraemer \& Ratamess, 2004). La ventaja práctica del modelo es
arquitectónica: demuestra la integración de inferencia local en el
dispositivo con respaldo por reglas.

\section{Discusión}
\label{sec:discusion}

\subsection{Frente al problema de adherencia}

La introducción parte del abandono temprano en los gimnasios. La evidencia
revisada indica que menos del 40\,\% de los socios mantiene un ejercicio regular
a largo plazo (Gjestvang et~al., 2021), por lo que las barreras de orientación,
estructura y seguimiento que abordan los objetivos específicos son pertinentes.
Sin embargo, este proyecto no midió si la aplicación mejora la adherencia: la
hipótesis queda sin contrastar y requiere un estudio con usuarios reales,
consentimiento informado y un periodo de seguimiento (sección~\ref{sec:recomendaciones}).

\subsection{Frente a la visión artificial para la ejecución de ejercicios}

Silva Amado (2024) propone un sistema de visión artificial que estima la pose
con MediaPipe, calcula ángulos articulares y los compara con umbrales para
identificar la ejecución de ejercicios, y lo evalúa de forma cualitativa. Los
enfoques son complementarios: aquel verifica \emph{cómo} se ejecuta un
movimiento; el presente trabajo reconoce \emph{con qué máquina} se entrena,
genera la rutina y registra el progreso, y aporta métricas cuantitativas
reproducibles. En consecuencia, la instrucción de la ejecución correcta (OE2)
se resuelve aquí con contenido guiado y no con verificación automática de la
postura, que se identifica como línea de trabajo futuro.

\subsection{Sobre la realidad aumentada}

El escáner combina lo real y lo virtual en todos los dispositivos, lo que
cumple la primera característica de la realidad aumentada descrita por Azuma
(1997). En los teléfonos compatibles, el módulo \archivo{CoachAR} añade el
registro en tres dimensiones: una vez reconocida la máquina, la ficha se fija a
un punto del entorno mediante un ancla de ARCore y se mantiene allí mientras la
persona camina o gira el teléfono (sección~\ref{subsec:it3}). Ese es un avance
frente a la versión anterior, que solo superponía la ficha en pantalla. Con todo,
la afirmación debe acotarse en tres sentidos. Primero, se comprobó su funcionamiento junto con el modelo de reconocimiento
en un dispositivo físico; esta prueba funcional no permite generalizar el
resultado a otros equipos ni cuantificar la precisión del anclaje.
Segundo, el ancla se fija a un punto de una superficie estimado a partir de la
caja del detector y no a la geometría de la máquina; no corrige errores de
profundidad ni sigue a la máquina si se desplaza. Tercero, el reconocimiento
sigue ocurriendo por captura y no de forma continua, así que la interacción en
tiempo real es parcial. Además, ARCore solo funciona en teléfonos certificados;
por eso la cámara convencional se mantiene como alternativa. Hasta contar con
mediciones (error de colocación, pérdidas de seguimiento, modelos de teléfono
probados), no hay evidencia de que la realidad aumentada mejore la experiencia
ni la adherencia; la revisión sistemática de Vinolo Gil et~al. (2021), en un
contexto distinto (fisioterapia), concluye que su eficacia es potencial pero no
concluyente.

\section{Amenazas a la validez y limitaciones}
\label{sec:amenazas}

\begin{itemize}
    \item \textbf{Validez de constructo del reconocedor.} Las 399 fotografías
    provienen de catálogos y no de gimnasios; el escenario de campo simulado
    aproxima, pero no reemplaza, las condiciones reales de luz, encuadre y
    oclusión.
    \item \textbf{Validez del detector.} Las métricas de la versión final
    (YOLO26m) se calcularon sobre la partición de validación, que también sirve
    para elegir el punto de control, y no sobre una partición de prueba
    independiente; la pérdida de validación, muy superior a la de entrenamiento,
    sugiere sobreajuste. 
    La primera versión se evaluó en solo 43 imágenes y 65 cajas, de las cuales 24 pertenecían a una sola clase; varias
    imágenes parecían series de la misma máquina desde ángulos cercanos. Ninguna
    versión se midió en un teléfono: no se conoce la latencia del \emph{endpoint},
    que escala a cero y puede tardar casi un minuto en responder tras un período
    inactivo (Ultralytics, s.~f.-b). El dispositivo de prueba es uno solo
    (Tabla~\ref{tab:dispositivo-prueba}), de gama media, y no se probaron otros
    equipos. El umbral de confianza de 0,8 se calibró
    para el modelo local y no se recalibró para el detector.
    \item \textbf{Validación de la realidad aumentada.} Se comprobó el funcionamiento
    integrado de ARCore y el modelo de reconocimiento en un dispositivo físico.
    No se realizaron mediciones del error de colocación ni de la latencia. Además, la detección de superficies depende de los rasgos visuales
    de la escena, por lo que suelos uniformes y espejos de un gimnasio podrían
    dificultarla (Google for Developers, s.~f.-b).
    \item \textbf{Dependencia de un servicio externo y su licencia.} El modo en
    la nube depende de la disponibilidad de Ultralytics Platform. La biblioteca
    Ultralytics se distribuye con licencia AGPL-3.0 o empresarial (Ultralytics,
    s.~f.-c), por lo que, si el producto se publica más allá del ámbito
    académico, habrá que revisar qué licencia corresponde.
    \item \textbf{Circularidad del recomendador.} Los datos sintéticos están
    etiquetados por las reglas que el modelo imita; la métrica no evalúa la
    utilidad para el usuario.
    \item \textbf{Ausencia de evaluación con usuarios.} No se midieron
    usabilidad, satisfacción ni adherencia. 
    \item \textbf{Cobertura de pruebas.} Las pruebas automatizadas cubren la
    lógica pura, la interpretación de la respuesta del reconocimiento, la
    geometría de la vista de ARCore y las reglas de seguridad; no cubren la
    interfaz ni se midieron la latencia y la exactitud de los modelos en un
    dispositivo físico.
    \item \textbf{Plataforma y conectividad.} El sistema funciona en Android y
    requiere conexión a internet; la persistencia sin conexión no está
    configurada.
    \item \textbf{Alcance nutricional.} El registro de alimentos se limita al
    conteo calórico y no sustituye la supervisión de un nutricionista.
\end{itemize}
